\documentclass[oneside]{scrbook}

% \usepackage[utf8]{inputenc}
\usepackage[T1]{fontenc}
\usepackage{version}
% \usepackage[latin1]{inputenc}
\usepackage[italian]{babel}
\usepackage{algpseudocode}
\usepackage{amsmath}
\usepackage{booktabs}
\usepackage{color}
\usepackage{amsfonts}
\usepackage{url}
\usepackage{algorithm}
\usepackage{chngcntr}
\usepackage[answerdelayed]{exercise}
\usepackage{graphicx}
\usepackage{caption}
\usepackage{listings}
\usepackage{tikz}
\usepackage{pgf}
\usepackage{colortbl}
\usepackage{listingsutf8}
\usepackage[pdftitle = {Dispense di esercizi}, pdfauthor = {Gianmarco Tedeschi e Alberto Marchesi}, pdfsubject = {Fondamenti di Informatica ($082746$ - A.A.~$2025-2026$)}, pdfcreator = {Gianmarco Tedeschi e Alberto Marchesi}, pdfproducer = {-}]{hyperref}
% \usepackage[automark,nouppercase]{scrpage2}
\usepackage{mdframed}
\usepackage{scrlayer-scrpage} % Necessary package for page styles

%\pagestyle{scrheadings}

%\usepackage{titlesec}
%\newcommand{\sectionbreak}{\clearpage}
\counterwithin{Exercise}{chapter}
\counterwithin{Answer}{chapter}

\algrenewcommand\algorithmicthen{\textbf{Allora}}
\algrenewcommand\algorithmicif{\textbf{Se}}
\algrenewcommand\algorithmicelse{\textbf{Altrimenti}}
\algrenewcommand\algorithmicend{\textbf{Chiudi}}
\algrenewcommand\algorithmicwhile{\textbf{Finch\'e}}
\algrenewcommand\algorithmicdo{\textbf{Esegui}}
\algrenewcommand{\algorithmiccomment}[1]{\hskip3em\texttt{//#1}}

\setlength{\parindent}{0in}
\setlength{\parskip}{10pt}

\newmdenv[%
        outerlinewidth=2,%
        roundcorner=10pt,%
        leftmargin=0,%
        rightmargin=0,%
        outerlinecolor=black,
        innertopmargin=\topskip,%
        splittopskip=\topskip,%
]{algo}


\ohead{\headmark}
\ihead{}
\chead{}
\ifoot{A.A.~$2025-2026$}
\cfoot{Fondamenti di Informatica (082746) - G. Tedeschi}
\ofoot{Pag.~\pagemark}
\setheadsepline{.4pt}
\setfootsepline{.4pt}

\definecolor{lbcolor}{rgb}{0.9,0.9,0.9}

\lstdefinestyle{cstyle}{%
  inputencoding=utf8/latin1,
  breaklines=true,
  breakatwhitespace=false,
  basicstyle=\small\ttfamily,
  language=c,
  frame=single,
  showstringspaces=false
}

\lstdefinestyle{cstylesmall}{%
  inputencoding=utf8/latin1,
  breaklines=true,
  breakatwhitespace=false,
  basicstyle=\scriptsize\ttfamily,
  language=c,
  frame=single,
  showstringspaces=false
}

\lstdefinestyle{console}{%
  inputencoding=utf8/latin1,
  breaklines=true,
  breakatwhitespace=false,
  basicstyle=\scriptsize\ttfamily,
  showstringspaces=false,
  backgroundcolor=\color{lbcolor}
}


\lstnewenvironment{clang}{\lstset{style=cstyle}}{}
\lstnewenvironment{clangsmall}{\lstset{style=cstylesmall}}{}


\addtocounter{chapter}{4}

\begin{document}

\chapter{Stringhe e Matrici}

\paragraph{Stringhe}
Le stringhe di caratteri sono gestite in C come dei vettori di \texttt{char} con alla fine un ``tappo'' dato dal carattere \texttt{'\textbackslash0'}. 

E' possibile acquisire un'intera stringa di caratteri in una sola istruzione grazie all'istruzione \texttt{scanf(`\%s`, variabile)} oppure \texttt{gets(variabile)}. Mentre l'istruzione \texttt{scanf(`\%s`, variabile)} ferma l'acquisizione al primo spazio,  
\texttt{gets(variabile)} inserisce nella variable tutto l'input fino al primo invio. \`E importante notare come quando si usano funzioni come \texttt{scanf()} o \texttt{gets()} per leggere input, se si immettono più caratteri o caratteri errati (ad esempio, superando i limiti di una stringa o inserendo caratteri non validi per un determinato tipo di input), questi caratteri possono rimanere nel buffer di input. Questo può creare problemi nelle successive richieste di input, perché queste potrebbero leggere direttamente dal buffer rimanente invece di attendere nuovo input dall'utente. Alcune soluzioni prevedono di \emph{svuotare} i caratteri rimanenti in modo da avere il buffer libero per la prossima lettura. Questo può essere eseguito utilizzando la funzione \texttt{getchar()} che consuma un singolo carattere, oppure utilizzando la funzione \texttt{fflush(stdin)} che svuoterà l'intero buffer di input.



Poich\`e le stringhe necessitano di un tappo, ogni volta che ne dichiaro una devo aggiungere un elemento per il tappo, ossia se voglio una stringa lunga \texttt{N\_MAX} dovr\`o dichiarare un array di \texttt{N\_MAX+1} elementi.

Esiste una libreria di C che gestisce le stringhe \texttt{string.h} essa ci permette di:
\begin{itemize}
		\item ottenere la lunghezza di una stringa con la funzione \texttt{strlen(stringa)};
		\item confrontare due stringhe \texttt{strcmp(stringa1, stringa2)} e restituisce zero se sono uguali, un numero negativo se \texttt{stringa1} viene prima in ordine lessicografico di \texttt{stringa2} e un numero positivo se \texttt{stringa2} viene prima in ordine lessicografico di \texttt{stringa1} (se i caratteri sono tutti maiuscoli o tutti minuscoli);
		\item copiare una stringa \texttt{stringa2} in un'altra stringa \texttt{stringa1} con la funzione\\
		\texttt{strcpy(stringa1, stringa2)};
		\item concatenare alla stringa \texttt{stringa1} la stringa \texttt{stringa2} con la funzione\\
		\texttt{strcat(stringa1,stringa2)};
\end{itemize}

\paragraph{Matrici}
Le matrici in C vengono definite come vettori di vettori. Ad esempio per dichiarare una matrice bidimensionale dovremo scrivere:
\begin{lstlisting}[style=cstyle]
#define DIM1 n_righe 
#define DIM2 n_colonne

void main(){

tipo_variabile nome_variabile[DIM1][DIM2];

istruzioni;
}
\end{lstlisting}
dove \texttt{n\_righe} e \texttt{n\_colonne} sono degli interi che specificano il numero di righe e colonne della matrice. Similmente ai vettori le matrici vengono memorizzate riga per riga, per indice di riga crescente, e, all'interno di ogni riga, per indice di colonna crescente (Figura \ref{fig:matrice}).

\begin{figure}[h!]
    \centering
    \begin{tikzpicture}
        % Draw the matrix grid (6x1)
        \draw[thick] (0,0) rectangle (6,-6); % Outer rectangle
        \foreach \i in {1,2,3,4,5} {
            \draw[thick] (0,-\i) -- (6,-\i); % Horizontal lines
        }

        % Matrix elements (texts)
        \node at (3,-0.5) {\texttt{val.\ matrice[0][0]}};
        \node at (3,-1.5) {\texttt{val.\ matrice[0][1]}};
        \node at (3,-2.5) {\texttt{val.\ matrice[0][2]}};
        \node at (3,-3.5) {\texttt{val.\ matrice[1][0]}};
        \node at (3,-4.5) {\texttt{val.\ matrice[1][1]}};
        \node at (3,-5.5) {\texttt{val.\ matrice[1][2]}};

        % Draw the red curly brace to group the first three cells for matrice[0]
        \draw[thick, red] (0,-0.1) -- (-0.5,-0.1) .. controls (-0.5, -1.5) and (-0.5, -1.5) .. (-0.5,-2.9) -- (0,-2.9);
        \node[red] at (-1.6, -1.5) {\texttt{matrice[0]}};

        % Draw the red curly brace to group the first three cells for matrice[1]
        \draw[thick, red] (0,-3.1) -- (-0.5,-3.1) .. controls (-0.5, -4.5) and (-0.5, -4.5) .. (-0.5,-5.9) -- (0,-5.9);
        \node[red] at (-1.6, -4.5) {\texttt{matrice[1]}};
        
    \end{tikzpicture}
    \caption{Rappresentazione di una matrice 2x3 in memoria}
    \label{fig:matrice}
\end{figure}

Nel caso di dimensioni superiori a $2$, dovremo mettere tante parentesi quadre e specificare le grandezze massime per ogni dimensione.
		
\section{Esercizi}

% stringhe
\begin{Exercise}[label={ex:strcat}]
  Date due stringhe immesse dall'utente, implementare le funzioni \texttt{strlen()}, \texttt{strcmp()}, \texttt{strcpy()} e \texttt{strcat()} senza utilizzare le funzioni di libreria \texttt{string.h}. Considerare delle stringhe in ingresso di al massimo $20$ caratteri.
\end{Exercise}

\begin{Answer}[ref={ex:strcat}]
  \lstinputlisting[style=cstylesmall]{src/my_strlen.c}
  \lstinputlisting[style=cstylesmall]{src/my_strcmp.c}
  \lstinputlisting[style=cstylesmall]{src/my_strcpy.c}
  \lstinputlisting[style=cstylesmall]{src/my_strcat.c}
\end{Answer}

\begin{Exercise}[label={ex:occorrenze}]
    Si scriva un programma che, lette da tastiera due stringhe, calcoli il numero di occorrenze della prima stringa nella seconda.
\end{Exercise}

\begin{Answer}[ref={ex:occorrenze}]
  	\lstinputlisting[style=cstylesmall]{src/occorrenze.c}
\end{Answer}

\begin{Exercise}[label={ex:palindromo}]
		\Question Scrivere un programma che determini se una frase acquisita da tastiera \`e palindroma. Controllare se sono palindrome le seguenti frasi (senza spazi e maiuscole):
		\begin{itemize}
				\item Eri un nano non annuire
				\item Ad una vera pia donna dei simili fili misi e annodai: pareva nuda
				\item O mordo tua nuora, o ari un autodromo
				\item Occorre portar aratro per Rocco
		\end{itemize}
		\Question (Bonus) Modificare il programma precedente per considerare stringhe in input contenenti anche caratteri speciali e spazi.
\end{Exercise}

\begin{Answer}[ref={ex:palindromo}]
		\lstinputlisting[style=cstylesmall]{src/palindroma.c}
\end{Answer}
\begin{Exercise}[label={ex:secret}]
    Si scriva un programma che legga da tastiera due stringhe, una denominata \texttt{msg} ed una \texttt{secret} tutte di caratteri minuscoli.
    Si vuole sapere se tutti i caratteri di secret sono contenuti, nell'ordine, in \texttt{msg}.
    In caso positivo, il programma deve stampare "ACCESSO CONSENTITO" e stampare una copia di \texttt{msg}, dove però ad ogni lettera riconosciuta come parte di \texttt{secret} viene
    sostituita la lettera maiuscola.
\end{Exercise}

\begin{Answer}[ref={ex:secret}]
    \lstinputlisting[style=cstylesmall]{src/secret.c}
\end{Answer}

\begin{Exercise}[label={ex:cesare}]
		Implementare la seguente variante del cifrario di Cesare\footnote{\url{http://it.wikipedia.org/wiki/Cifrario_di_Cesare}}. Dopo aver acquisito una messaggio di massimo $160$ caratteri (alfabetici minuscoli), il programma dovr\`a chiedere all'utente una chiave (numero intero $K$, $1 < K < 25$).

  	Il programma stamper\`a il messaggio cifrato, ottenuto traslando ogni lettera di $K$ posizioni in avanti (dove $K$ \`e la chiave).

  	Dopo aver stampato il messaggio cifrato, il programma chiede all'utente di inserire un messaggio (che si assume sia stato cifrato con lo stesso algoritmo di cui sopra). Tale messaggio sar\`a dunque decifrato dal programma, il quale dovr\`a svolgere l'operazione inversa.

  	Infine, il messaggio decifrato sar\`a stampato a video.
		
		Suggerimento: i caratteri minuscoli dalla a alla z corrispondono nella tabella ASCII alle cifre intere dalla $97$ alla $122$.
\end{Exercise}

\begin{Answer}[ref={ex:cesare}]
  \lstinputlisting[style=cstylesmall]{src/cesare.c}
\end{Answer}
% \begin{Exercise}[label={ex:farfallino}]
  \Question Scrivere un programma che converta una stringa in alfabeto farfallino. Nell'alfabeto farfallino, ogni vocale \`e raddoppiata, e tra le due vocali \`e inserita la lettera `f'. Ad esempio, ``ciao'' diventa ``cifiafaofo''.

  \Question (bonus) Scrivere il medesimo programma utilizzando soltanto una variabile per memorizzare la stringa tradotta (ovvero senza copiare la traduzione in una nuova variabile).
\end{Exercise}

\begin{Answer}[ref={ex:farfallino}]
  \Question
  \lstinputlisting[style=cstylesmall]{src/farfallino.c}

  \Question
  \lstinputlisting[style=cstylesmall]{src/farfallino-bonus.c}
\end{Answer}
\begin{Exercise}[label={ex:istogramma}]
  Scrivere un programma che considera solo le lettere maiuscole e minuscole di una stringa letta da tastiera (di dimensione massima $256$ elementi). Il programma dovr\`a calcolare le occorrenze di ogni carattere nella stringa. Le occorrenze saranno poi stampate a video a video come istogramma con gli asterischi nel seguente modo:

\begin{lstlisting}[style=console]
str = Ciao Come Stai? Non c'c' male GRAZIE! Mi trovo molto bene in questa citta'.

a |
b | *
c | **
d |
e | ******
f |
g |
h |
i | *****
j |
k |
l | **
m | ***
n | ***
o | *******
p |
q | *
r | *
s | *
t | ******
u | *
v | *
w |
x |
y |
A |
B |
C | **
D |
E | *
F |
G | *
H |
I | *
J |
K |
L |
M | *
N | *
O |
P |
Q |
R | *
S | *
T |
U |
V |
W |
X |
Y |
\end{lstlisting}
\end{Exercise}

\begin{Answer}[ref={ex:istogramma}]
  \lstinputlisting[style=cstylesmall]{src/istogramma.c}
\end{Answer}
\begin{Exercise}[label={ex:anagramma}]
  Scrivere un programma che determini se due parole acquisite da tastiera sono una l'anagramma dell'altra.
\end{Exercise}


\begin{Answer}[ref={ex:anagramma}]
		\lstinputlisting[style=cstylesmall]{src/anagramma_2.c}
        \paragraph{Nota: } Possiamo risolvere lo stesso algoritmo senza stringhe di appoggio ed in maniera più efficiente.
		\lstinputlisting[style=cstylesmall]{src/anagramma.c}
\end{Answer}

% matrici
\begin{Exercise}[label={ex:matrice_diagonale}]
		Scrivere un programma che chieda all'utente la dimensione effettiva di una matrice quadrata, che dovrà poi essere acquisita dal programma.

  	Successivamente, il programma dovrà controllare se la matrice è diagonale (i.e., nessuno zero sulla diagonale principale e solo zeri fuori da essa).
  	
  	Estendere il programma per controllare anche se la matrice inserita è simmetrica.
  	
\end{Exercise}

\begin{Answer}[ref={ex:matrice_diagonale}]
  	\lstinputlisting[style=cstylesmall]{src/matrice_diagonale.c}
\end{Answer}

\begin{Exercise}[label={ex:spirale}]
  Scrivere un programma che inizializzi una matrice $10 \times 10$ con i valori della tavola pitagorica.

  Dopodich\'e, il programma dovr\`a stampare la matrice seguendo un percorso a spirale in senso orario, a partire dalla cella $m[0][0]$.

  Per semplicit\`a si pu\`o assumere la matrice quadrata.
\end{Exercise}

\begin{Answer}[ref={ex:spirale}]
  \lstinputlisting[style=cstylesmall]{src/matrice-spirale.c}
\end{Answer}

\begin{Exercise}[label={ex:istogramma_matrice}]
  	Data $m$ una matrice di dimensione $n \times n$ (costante simbolica) di numeri interi nell'intervallo $[0,9]$:
  	\begin{enumerate}
  			\item si scriva un frammento di programma in linguaggio C (con le relative dichiarazioni di variabili necessarie al corretto funzionamento), che trovi il numero pi\`u frequente (si ipotizzi che tale numero sia unico).
				\item si scriva un frammento che memorizzi in un vettore e stampi a schermo tutti i numeri presenti nella matrice con valore minore del valore pi\`u frequente e successivamente memorizzi in un altro vettore tutti quelli con valore maggiore di quello più frequente e li stampi;
  			\item aggiungere al programma di cui sopra un frammento di codice che legga tutti e soli i valori memorizzati nel secondo dei due vettori e stampi a video se i valori che vi sono contenuti sono o meno monotoni crescenti, ovvero se per ogni $i$ si ha $a_i \leq a_{i+1}$.
  	\end{enumerate}
\end{Exercise}

\begin{Answer}[ref={ex:istogramma_matrice}]
  	\lstinputlisting[style=cstylesmall]{src/opmatrice.c}
\end{Answer}
\begin{Exercise}[label={ex:riga_somma_massima}]
		Scrivere un programma che acquisisca una matrice di interi positivi con \texttt{DIM\_R} righe e \texttt{DIM\_C} colonne. Il prgramma deve stampare l'indice della riga della matrice con somma massima.
  	
\end{Exercise}

\begin{Answer}[ref={ex:riga_somma_massima}]
  	\lstinputlisting[style=cstylesmall]{src/riga_somma_massima.c}
\end{Answer}

\begin{Exercise}[label={ex:temperatura}]
    Una stazione meteorologica registra la temperatura esterna ogni ora, ogni giorno per un mese intero, in una struttura dati apposita: una matrice bi-dimensionale di dimensioni fisiche di 31 righe (i giorni del mese) per 24 colonne (le ore del giorno).
    Si scriva un programma che calcoli i seguenti dati:
    \begin{itemize}
        \item Temperatura media del mese.
        \item Temperatura media giornaliera.
        \item Temperatura media diurna (ore comprese tra le 7:00 e le 19:00), su tutto il mese.
        \item Temperatura media notturna su tutto il mese.
    \end{itemize}

    \textbf{Nota: } non tutti i mesi hanno 31 giorni! Popoliamo la matrice utilizzando la funzione \texttt{rand()} della libreria \texttt{stdlib}.
\end{Exercise}

\begin{Answer}[ref={ex:temperatura}]
  	\lstinputlisting[style=cstylesmall]{src/temperatura.c}
\end{Answer}


\clearpage
\section*{Soluzioni}
\shipoutAnswer

\end{document}